\documentclass[../main.tex]{subfiles}
\begin{document}

% --------------------------------------------------
\section{theorem.sty}
% --------------------------------------------------

\texttt{theorem.sty}では, 定理, 補題, 命題, 定義, 例, 注意などの環境を定義している. 
見た目の作成には\texttt{tcolorbox}を使い, 色は\texttt{core.sty}で定義した\verb|\ThmColor|や\verb|\DefiColor|などを参照している. 

\subsection{定理系環境}

用意している定理系環境は以下の通りである. 
\begin{itemize}
	\item \verb|thm| : 定理
	\item \verb|lem| : 補題
	\item \verb|prop| : 命題
	\item \verb|cor| : 系
	\item \verb|defi| : 定義
	\item \verb|ex| : 例
	\item \verb|rem| : 注意
\end{itemize}

これらの環境は, いずれも
\begin{center}
	\verb|\begin{<環境名>}{<説明>}{<ラベル>} ... \end{<環境名>}|
\end{center}
の形で使う. 
\verb|<説明>|には定理名, 引用元, 簡単な補足などを書くことができる. 
\verb|<ラベル>|に指定した名前は, 後から\verb|\zcref|で参照できる. 

番号は節ごとにリセットされ, 定理, 補題, 命題, 系, 定義, 例, 注意で同じカウンタを共有する. 
したがって, 本文中では「定理1.1」の次に「定義1.2」が続くような番号付けになる. 

\begin{defi}{群}{doc-def-group}
	集合$G$と写像$G\times G\to G$が群であるとは, 結合法則, 単位元の存在, 逆元の存在を満たすことである. 
\end{defi}

\begin{thm}{準同型定理}{doc-thm-hom}
	$f:G\to H$を群準同型とする. 
	このとき, $f$から誘導される写像
	\begin{align}
		\tilde{f}:G/\Ker f\to \Im f,\quad x\Ker f\mapsto f(x) \label{doc-thm-hom-map}
	\end{align}
	は同型写像である. 
\end{thm}

\begin{proof}
	$x\Ker f=y\Ker f$ならば$xy^{-1}\in\Ker f$である. 
	したがって$f(x)=f(y)$となり, \zcref{doc-thm-hom-map}はwell-definedである. 
	全射性は像の定義から従い, 単射性も同様に核の定義から従う. 
	よって\zcref{doc-thm-hom}が示された. 
\end{proof}

\subsection{証明環境}

\texttt{amsthm}の\texttt{proof}環境は, \texttt{tcolorbox}で囲む形に調整している. 
証明記号は\verb|\AccentColor|の黒四角になる. 

証明の見出しを変えたい場合は, 通常の\texttt{proof}環境と同じく
\begin{center}
	\verb|\begin{proof}[別証] ... \end{proof}|
\end{center}
のように書けばよい. 

\subsection{Pythonコード}

現在, Pythonコードを掲載するための\verb|pycode|環境も\texttt{theorem.sty}で定義している. 
これはSageMathのコードをノートに載せることを想定した環境であり, \texttt{tcolorbox}と\texttt{listings}を使っている. 

\begin{center}
	\verb|\begin{pycode}{<タイトル>}{<ラベル>} ... \end{pycode}|
\end{center}
の形で使う. 
ラベルの型は\texttt{listing}なので, \verb|\zcref|で参照すると「コード」と表示される. 

\begin{pycode}{繰り返し二乗法}{doc-py-power}
def power(base, exp):
	result = 1
	while exp != 0:
		if exp & 1 == 1:
			result *= base
		base *= base
		exp >>= 1
	return result
\end{pycode}

例えば, \zcref{doc-py-power}のように本文中から参照できる. 

\end{document}
